% ===================================================================
%  TABULKA č. 16 — Velký obchodní dům. — Bazar. — Hračky.
%  Traduction 2026 d'apres texto/io, controlee sur texto/fr, texto/ru
%  et texto/uk. Consigne : voir texto/cs/10-tabulka-01.tex.
%
%  Le tableau d'astronomie (bloc 15-1) porte des renvois a LETTRES,
%  (a) a (f), graves sur la planche elle-meme : ils se rangent sous
%  le numero 85, comme le dit gravuri/literi.json. Ce bloc est aussi
%  le troisieme ecart declare de renvoji.py.
%
%  LA MENAGERIE DE CARTON DONNE LE MEME RESULTAT QU'EN IRLANDAIS ET
%  EN GALICIEN : les betes qu'on a vues gardent leur mot —
%  « liška » le renard, « vlk » le loup, « myš » la souris,
%  « netopýr » la chauve-souris, « lasice » la belette, « veverka »
%  l'ecureuil, « kos » le merle, « slavík » le rossignol — et les
%  betes qu'on n'a jamais vues arrivent avec le livre d'images :
%  « hroch », « dromedár », « leopard », « hyena », « krokodýl »,
%  « panter », « pštros ».
%
%  ET LE TCHEQUE AJOUTE ICI CE QUE NI L'IRLANDAIS NI LE GALICIEN NE
%  POUVAIENT DONNER. Deux de ces betes lointaines y portent quand
%  meme un nom DU PAYS : « nosorožec », le porte-corne, pour le
%  rhinoceros, et « velbloud » pour le chameau. La regle du tableau
%  11 se lit donc jusque dans une boutique de jouets : ce qui n'est
%  jamais venu peut recevoir un nom du pays, si quelqu'un decide de
%  le lui donner.
%
%  DERNIER TABLEAU DE LA COLONNE TCHEQUE.
% ===================================================================

%%K t16-apar-1 apar
\VUsaut{4.0mm}
\VUtitre{15.0pt}{60}{TABULKA č. 16}
\VUsaut{2.0mm}
\VUpk{13.2pt}{40}{Velký obchodní dům. --- Bazar.}
\VUpk{13.2pt}{40}{Hračky.}
\VUsaut{3.4mm}

%%K t16-01-1 p
1. --- Nový rok se blíží. Velké obchodní domy připravily své výklady \VUgras{hraček}
\textsuperscript{(1)} a novoročních dárků.

%%K t16-01-2 p
Požádal jsem svého dědečka, aby mě zavedl do bazaru, abych viděl tu krásnou výstavu, a
on svolil.

%%K t16-02-1 p
2. --- Nejprve jsme se zastavili před krásnými zbrojními štíty, které obsahují
\VUgras{pušku,} \VUgras{čapku,} \VUgras{šavli,} \VUgras{opasek na kord} a pár
\VUgras{nárameníků,} nebo \VUgras{přilbu,} \VUgras{pancíř} \textsuperscript{(3)} a
\VUgras{pistoli} \textsuperscript{(4)}. To vše mě zajímalo mnohem více než
\VUgras{světlomety} \textsuperscript{(2)}.

%%K t16-tit-1 sub
\VUsaut{3.0mm}
\VUpk{10.2pt}{40}{Krásné podívané.}
\VUsaut{2.6mm}

%%K t16-03-1 p
3. --- V témž oddělení byl \VUgras{strážný} \textsuperscript{(5)} ve své
\VUgras{budce}; zdálo se, že stojí na stráži před celou lepenkovou zvěřincí. Kolik
zvířat tam bylo: \VUgras{papoušek} \textsuperscript{(6)} na svém bidélku,
\VUgras{nosorožec} \textsuperscript{(7)} s rohem na nose, \VUgras{sob}
\textsuperscript{(8)}, \VUgras{pštros} \textsuperscript{(9)}, \VUgras{opice}
\textsuperscript{(10)} louskající ořech, \VUgras{liška} \textsuperscript{(11)}
odnášející slepici, \VUgras{krokodýl} \textsuperscript{(12)} zející obrovskými
čelistmi, \VUgras{velbloud} \textsuperscript{(13)}, elegantní \VUgras{chrt}
\textsuperscript{(14)}, \VUgras{panter} \textsuperscript{(15)}, potom dvě zvířata,
která jsem neznal a která jsou, jak se říká, \VUgras{lasice} \textsuperscript{(16)} a
\VUgras{hyena} \textsuperscript{(17)}. Také tam byl \VUgras{leopard}
\textsuperscript{(18)} a \VUgras{vlk} \textsuperscript{(21)}; zcela blízko byly
roztomilé \VUgras{krysičky} \textsuperscript{(19)} a maličké gumové \VUgras{myšky}
\textsuperscript{(20)}. Nakonec \VUgras{hroch} \textsuperscript{(22)} a hrbatý
\VUgras{dromedár} \textsuperscript{(23)} uzavírali sbírku. Zůstal bych tam mnoho hodin,
kdyby nebylo blízko skvělého tábora plného \VUgras{olověných vojáčků}
\textsuperscript{(29)}, který upoutal mou pozornost.

%%K t16-tit-2 sub
\VUsaut{4.0mm}
\VUpk{10.2pt}{40}{Vojáci a válka.}
\VUsaut{2.8mm}

%%K t16-04-1 p
4. --- Můj dědeček, který je \VUgras{invalida} \textsuperscript{(24)} a který musel
opustit armádu pro \VUgras{zranění,} jež mu vyneslo \VUgras{vojenskou medaili,} má
velmi rád, když mi vypravuje o \VUgras{taženích,} jichž se zúčastnil. Byl šťasten,
když vysvětloval různé pohyby malé armády v miniatuře. Skupina opravdových vojáků, kteří
navštívili bazar --- \VUgras{ženista} \textsuperscript{(25)}, \VUgras{alpský myslivec}
\textsuperscript{(26)} s hlavou pokrytou \VUgras{baretem,} \VUgras{štábní šikovatel}
\textsuperscript{(27)} pěchoty a \VUgras{dělostřelecký šikovatel}
\textsuperscript{(28)} se zlatým \VUgras{prýmkem} na rukávě --- se zastavila blízko
nás, aby vyslechla vysvětlení.

%%K t16-05-1 p
5. --- Můj dědeček, zcela hrdý na to, že měl tak skvělé posluchače, improvizoval úplný
plán bitvy.

%%K t16-05-2 p
Opevněné město se svými \VUgras{hradbami} a \VUgras{příkopy} bylo \VUgras{obléháno}
\VUgras{nepřátelskou armádou}, která byla po dlouhém \VUgras{bombardování} připravena k
\VUgras{útoku}. Ale \VUgras{obležení}, donuceni hladem, se pokusili o \VUgras{výpad}
proti \VUgras{táboru} \VUgras{obléhatelů}. Když odrazili \VUgras{útok} úporným
\VUgras{bojem} kolem \VUgras{tábora}, vyslali \VUgras{vítězní} obléhatelé své
\VUgras{eskadrony}, aby pronásledovaly \VUgras{poražené} útočníky. Ti
\VUgras{ustoupili} pokrývajíce \VUgras{bojiště} \VUgras{mrtvými} a \VUgras{raněnými}.
\VUgras{Nosiči nosítek} odnesli ty poslední do \VUgras{ambulance,} aby je ošetřil
\VUgras{chirurg}.

%%K t16-05-3 p
Přes hrdinný odpor několika \VUgras{praporů,} které vytvořily \VUgras{čtverec,} byla
\VUgras{porážka} úplná, zbytek armády \VUgras{prchl} a zanechal mnoho \VUgras{zajatců}.
Nepřítel šel dokonat své \VUgras{vítězství} obsazením města. Snad to bude konec tažení a
po \VUgras{příměří} bude možno podepsat \VUgras{mírovou smlouvu}.

%%K t16-tit-3 sub
\VUsaut{4.4mm}
\VUpk{10.2pt}{40}{Novoroční dárky.}
\VUsaut{2.8mm}

%%K t16-06-1 p
6. --- Mladí vojáci, kteří se smáli malebnému vypravování veterána, ho požádali o
několik dalších vysvětlení. Co se mne týče, obešel jsem obchod, abych se podíval na
krásnou \VUgras{kuši} \textsuperscript{(32)} a na \VUgras{luk} \textsuperscript{(33)} s
jeho \VUgras{šípem} \textsuperscript{(31)}. Tam jsem našel jinou sbírku zvířat: dva
\VUgras{zpěvné ptáky} \textsuperscript{(34)} --- automatického \VUgras{slavíka} a
\VUgras{pěnici} zpívající z plného hrdla --- a řadu podivných zvířat:
\VUgras{netopýra} \textsuperscript{(35)}, \VUgras{hroznýše} \textsuperscript{(36)},
\VUgras{želvu} \textsuperscript{(37)}, \VUgras{tuleně} \textsuperscript{(38)},
\VUgras{velrybu} \textsuperscript{(39)} a plyšového \VUgras{žraloka}
\textsuperscript{(40)}.

%%K t16-07-1 p
7. --- Nezastavil jsem se dlouho, abych je pozoroval, neboť jsem zahlédl to, o čem jsem
snil: překrásné \VUgras{loutkové divadlo} \textsuperscript{(41)} se skvělými
\VUgras{dekoracemi} a půvabnou červenou \VUgras{oponou.} Na \VUgras{jevišti} se zdálo,
že dva herci hrají \VUgras{roli} ve velmi zajímavé \VUgras{hře,} neboť malí lepenkoví
\VUgras{diváci} umístění v ložích nebo sedící v \VUgras{křeslech} a v \VUgras{přízemí}
za \VUgras{dirigentem} živě tleskali.

%%K t16-07-2 p
Bohužel to divadlo stálo velmi draho.

%%K t16-08-1 p
8. --- Ostatně vybrat hračky bylo nesnadné, neboť ve výkladech se hromadily hračky
všeho druhu: \VUgras{Harlekýn} \textsuperscript{(42)}, \VUgras{Kašpárek}
\textsuperscript{(43)}, \VUgras{Pierot} \textsuperscript{(44)}, \VUgras{výři}
\textsuperscript{(45)} mávající křídly, hvízdající \VUgras{kosi}
\textsuperscript{(46)}, štěkající \VUgras{pudlové}, \VUgras{fonograf}
\textsuperscript{(47)} a \VUgras{trpělivostní hry.} Jedna z těch hraček mě velmi
pobavila: představovala \VUgras{námořní bitvu} \textsuperscript{(48)}, v níž byla
\VUgras{loď} napadena \VUgras{piráty} blízko země osázené \VUgras{kokosovníky}.

%%K t16-noto-1 noto
\VUnotes{143.2mm}{%
\textit{(*) Viz tabulku č. 6.}%
}

%%K t16-09-1 p
9. --- Mohl jsem velmi snadno pozorovat všechny ty divy, neboť v obchodě bylo jen málo
lidí, sotva několik zevlounů, \VUgras{dragoun} \textsuperscript{(49)} a
\VUgras{inkasista} \textsuperscript{(50)}, který předkládal \VUgras{směnku} u hlavní
\VUgras{pokladny} \textsuperscript{(51)}.

%%K t16-10-1 p
10. --- Zeptal jsem se svého dědečka, k čemu slouží knihy položené na policích za
\VUgras{pokladní}; odpověděl mi, že jsou to \VUgras{účetní knihy}.

%%K t16-tit-4 sub
\VUsaut{3.6mm}
\VUpk{10.2pt}{40}{Zevlování kolem krámu.}
\VUsaut{2.8mm}

%%K t16-11-1 p
11. --- Když jsme opustili oddělení hraček, šli jsme do sedlářství podívat se na
\VUgras{sedla} \textsuperscript{(53)}, \VUgras{třmeny} \textsuperscript{(54)},
\VUgras{uzdy} \textsuperscript{(55)}, \VUgras{udidla} \textsuperscript{(56)} a
\VUgras{jezdecké biče.} Zatímco \VUgras{vedoucí oddělení} \textsuperscript{(57)} dával
mému dědečkovi několik informací, šel jsem se podívat na výklad \VUgras{porcelánu} a
\VUgras{křišťálu} \textsuperscript{(58)}, kde jsou krásné \VUgras{salátové mísy} a
jemné \VUgras{čajové konvice} \textsuperscript{(59)}.

%%K t16-12-1 p
12. --- Abychom šli do prvního patra, prošli jsme za výkladem \VUgras{korzetů}
\textsuperscript{(60)}, blízko obchodu s punčochami, kde byly vystaveny velmi krásné
\VUgras{kalhoty} \textsuperscript{(63)}. Blízko toho \VUgras{prodavačka}
\textsuperscript{(61)} dávala \VUgras{kupující} \textsuperscript{(62)} vybírat krásné
hedvábné \VUgras{látky} \textsuperscript{(64)}.

%%K t16-12-2 p
Když jsme vystupovali po schodišti, všiml jsem si, že obchod je vyzdoben japonskými
\VUgras{lampiony} \textsuperscript{(65)}, \VUgras{slunečníky} \textsuperscript{(66)} a
obrovskými \VUgras{palmami} \textsuperscript{(70)}.

%%K t16-13-1 p
13 . --- V prvním patře nebylo mnoho k vidění; jen nábytek (*), \VUgras{železné
pokladny} \textsuperscript{(67)}, \VUgras{komody} \textsuperscript{(68)} a
\VUgras{dětské kočárky} \textsuperscript{(69)}. Ale celkový pohled na obchod byl velmi
\VUgras{zábavný.}

%%K t16-14-1 p
14. --- U našich nohou jsem viděl hudební nástroje: \VUgras{harfy}
\textsuperscript{(71)}, \VUgras{kytary} \textsuperscript{(72)}, \VUgras{mandolíny}
\textsuperscript{(73)}, \VUgras{křídlovky} \textsuperscript{(74)},
\VUgras{klarinety} \textsuperscript{(75)}, housle s jejich \VUgras{smyčcem}
\textsuperscript{(76)} a dokonce \VUgras{velký buben} \textsuperscript{(77)}. Poněkud
dále, u papírnického oddělení, \VUgras{student} \textsuperscript{(78)} smlouval s
\VUgras{prodavačem} \textsuperscript{(79)} o desky na sešity. Blízko nich jsem rozeznal
krásný \VUgras{slabikář} \textsuperscript{(80)}.

%%K t16-14-2 p
Uprostřed obchodu byl vztyčen vysoký \VUgras{vánoční stromek}
\textsuperscript{(81)} zcela opatřený svíčkami, drobnostmi a všelikými hračkami:
\VUgras{dřevěnými koníky} \textsuperscript{(82)}, \VUgras{panenkami}
\textsuperscript{(83)} atd.

%%K t16-14-3 p
\VUgras{Voňavkářství} \textsuperscript{(84)} se svými \VUgras{zubními pastami} a
různými \VUgras{voňavkami} šířilo velmi dobrou vůni, která k nám doléhala.

%%K t16-tit-5 sub
\VUsaut{3.4mm}
\VUpk{10.2pt}{40}{Něco kupujeme.}
\VUsaut{2.8mm}

%%K t16-15-1 p
15. --- Protože můj dědeček začal shledávat čas příliš dlouhým, sešli jsme po velkém
schodišti. Chvíli se zastavil před \VUgras{astronomickou tabulkou}
\textsuperscript{(85)} znázorňující, jak mi řekl, \VUgras{komety}
\textsuperscript{(a)}, \VUgras{zatmění měsíce a slunce} \textsuperscript{(b)},
větrnou růžici se \VUgras{čtyřmi světovými stranami} \textsuperscript{(c)}
\VUgras{(sever, jih, východ a západ)}, mapu obou \VUgras{polokoulí}
\textsuperscript{(d)}, \VUgras{planety} \textsuperscript{(e)} a \VUgras{sluneční
soustavu} \textsuperscript{(f)}. Nerozuměl jsem z toho všeho ničemu a mnohem více mě
zajímala prýmkovaná livrej \VUgras{poslíčka} \textsuperscript{(86)}, který procházel, a
krásné \VUgras{vějíře} \textsuperscript{(87)}, které byly blízko mne, vedle
\VUgras{peněženek} \textsuperscript{(88)} a \VUgras{náprsních tašek}
\textsuperscript{(89)}.

%%K t16-16-1 p
16. --- Obdivoval jsem také na výkladě \VUgras{pera} \textsuperscript{(90)} a
\VUgras{přezky k opasku} \textsuperscript{(91)} a krásné \VUgras{umělé květiny}
\textsuperscript{(94)} půvabně uspořádané v koších. \VUgras{Fialky}, \VUgras{lilie},
\VUgras{hyacinty} \textsuperscript{(95)}, \VUgras{muškáty} \textsuperscript{(96)},
\VUgras{bílé lilie} \textsuperscript{(97)} a \VUgras{lichořeřišnice}
\textsuperscript{(98)} byly velmi dobře napodobeny.

%%K t16-tit-6 sub
\VUsaut{4.4mm}
\VUpk{10.2pt}{40}{Dárek od mého dědečka.}
\VUsaut{2.8mm}

%%K t16-17-1 p
17. --- Poprosil jsem svého dědečka, aby mi koupil jeden ze \VUgras{zubních kartáčků}
\textsuperscript{(92)}, které byly v krabici blízko \VUgras{sponek do vlasů}
\textsuperscript{(93)}. Svolil a já jsem šel sám zaplatit svůj nákup u pokladny, blízko
níž se připravovala významná \VUgras{zásilka} \textsuperscript{(100)} pro
\VUgras{zákazníka} \textsuperscript{(99)}.

%%K t16-18-1 p
18. --- Náhle nastal lehký zmatek mezi osobami, které se zastavily blízko uměleckých
\VUgras{bronzů} \textsuperscript{(101)}. \VUgras{Zloděj} \textsuperscript{(102)} byl
právě přistižen při činu \VUgras{inspektorem} \textsuperscript{(103)}, když se pokoušel
někoho okrást. Postarali se o to, aby přivedli strážníka a dali ho \VUgras{zatknout,} a
právě když jsme vycházeli, byl provinilec odváděn do vězení.
